% Paper 3 — Related Work (borrador 2026-10-17).
% Fuentes: docs/paper3-outline-v3-2026-10-17.md. Toda afirmación sobre trabajos ajenos se
% verificó contra su texto (HTML/PDF) o su abstract en arXiv en esta sesión o en las notas de
% lectura de agosto; las claves nuevas están en paper3/refs-related-2026-10-17.bib, las de
% agosto en docs/paper3-refs-verificadas-2026-08-28.bib y paper2/refs.bib.

\section{Related Work}\label{sec:related}

\paragraph{Automatic harness optimization and its acceptance rules.}
A line of work that began in early 2026 treats the agent harness, rather than
the model weights, as the object of optimization. Meta-Harness searches over
harness code with a proposer agent and keeps the candidate with the best score
on its search set \citep{lee2026metaharness}. AutoDesign formalizes the same
loop for long-horizon design tasks and accepts a candidate when its training
objective improves and its development objective does not degrade
\citep{luo2026autodesign}. Harness Continual Learning frames the harness as a
continually updated learning state and commits an update when it gains at
least $\delta$ on the current tasks while retaining an anchor set
\citep{kang2026hcl}. HarnessOpt-Bench separates search from evaluation with a
held-out partition that stays inaccessible throughout the search, and names
the optimizer's ability to ``separate real improvement from noise'' among the
capabilities it measures \citep{ursekar2026harnessopt}. All four decide what to
commit by comparing point estimates against a margin, and none of them derives
that margin from the run-to-run noise of the configuration being evaluated. The
false-acceptance rates in Section~\ref{sec:results} apply these decision rules,
not the systems that use them: we re-implement each rule on pairs of runs of an
identical configuration, where any accepted change is a false positive.

\paragraph{Statistical gates for self-evolving agents (concurrent work).}
Between June and September 2026, four systems replaced the point-estimate rule
with a statistical test. PACE accepts a prompt edit through a paired
testing-by-betting e-process over discordant instances, committing once the
wealth exceeds $1/\alpha$ with $\alpha$ fixed by the user; its guarantee is
explicitly per decision rather than per run, and its false commits are labelled
by a fresh audit pool rather than by ground truth \citep{zayxshawn2026pace}.
Its experiments are, notably, on small models (Qwen2.5 at 0.5--3B) and
arithmetic or multiple-choice tasks. HarnessBank gates offspring harnesses with
a paired $z$-test on per-task differences, $z \geq 1.96$, where the standard
deviation is computed across tasks within a single evaluation; it reports that
crediting a candidate from a single run, or from the mean of three, produces phantom progress in 62--76\% of post-convergence rounds with neutral candidates \citep{luo2026harnessbank}. SAGE applies a
one-sided paired test to skill edits and states that the significance level
``should match the task noise'', but tunes it on held-out data without
measuring that noise, and lists ``an adaptive threshold that tightens or
loosens $\alpha$ from an online estimate of task noise'' as future work
\citep{wang2026sage}. RRSI comes closest to our proposal: before evolution it
repeatedly evaluates the unchanged base harness and sets its noise tolerance $\delta$ from the spread of those evaluations (0.017, 0.004 and 0.020 on its three domains; on one of them, as the first percentile of a half-split null), then keeps $\delta$ fixed for the whole run, runs no sweep over
it, and does not report a false-acceptance rate \citep{xia2026rrsi}.

These gates establish that statistics belong in the accept step, and RRSI
establishes that the threshold can be calibrated against measured noise. What
none of them does is re-derive the threshold for each incumbent, or report how
the false-acceptance rate moves with the threshold. That is the gap this paper
fills. Our data show that the noise of a configuration is not a property of the
suite: within a single sweep, the share of tasks that flip between identical
runs differs across configurations by a median of 11.1 points (up to 26.3,
from 0\% to 26.3\% on one suite), and one incumbent in three is noisier than
its base. A tolerance measured once on the base therefore accepts
null changes at a rate its calibration does not control (about 7\%, and up to
18\% when the base happens to be the quietest configuration). Re-deriving it
per configuration lowers that rate below 1\% and matches or raises power at
every effect size we injected (Section~\ref{sec:results}).

\paragraph{Run-to-run noise in agent evaluation.}
A parallel literature measures the noise that these gates must contend with.
Recuris re-ran an unchanged package three days later, obtained an interval of
$[-6.98, +7.27]$, and concluded that differences of a few points should be
treated as run-to-run variation regardless of their interval
\citep{yu2026recuris}; it uses the floor as a qualitative caveat in one domain
rather than as the threshold of a rule. \citet{zhang2026stopcomparing}
decompose benchmark variance into model, harness and interaction terms; the
decomposition is exact for the cell means but omits the residual between runs
of the same cell, which their own two-run design estimates and does not report.
\citet{parupudi2026neutral} show that leaderboard orderings are manufactured
by configuration-fragile items under deterministic decoding, a source
orthogonal to ours (configuration rather than repetition) in a domain they
explicitly exclude from agentic tasks. More recently, the Noise Floor Audit finds reruns of three hosted tool-calling endpoints nearly deterministic at temperature 0 on the multiple and parallel categories of BFCL (0.7--2.7\% of items ever flip), while semantics-preserving prompt perturbations produce paired standard deviations 11 to 58 times larger \citep{chen2026noisefloor}; \citet{wiedmann2026systems} attribute about 54\% of
outcome variance on four scientific tasks to repeating the same configuration;
and \citet{hardy2026reliability} find model-scaffold rankings highly reliable
but model rankings much less so (0.148--0.841) across 22 benchmarks. On
SWE-bench Verified, exact McNemar tests separate none of the 29 adjacent pairs
in the top thirty, while the score of a single model varies by up to 29.8 points across scaffolds \citep{liu2026converged}. These studies measure the noise; our contribution is
to show what happens to an acceptance rule when that measurement is taken once
and reused.

\paragraph{Overfitting in harness evolution.}
Two benchmarks question the gains of harness evolution itself. HarnessDev
finds that evolution-stage revisions bring gains that ``are unstable and
transfer only partially to held-out tasks'' \citep{wu2026harnessdev}, and
\citet{wang2026rethinking} find that, on Terminal-Bench 2.1 and under matched
budgets, automatic harness evolution does not beat simple test-time scaling,
with long-horizon games as the exception. Both point to the accept step as a
place where noise enters the search. A gate whose false-acceptance rate is
uncontrolled commits candidates that will not replicate, and the instability
these benchmarks report is the aggregate symptom of that per-decision error.

\paragraph{Scope.}
Our regime differs from most of the work above: small local models (3--20B)
with large capability gaps, Rust coding suites graded by \texttt{cargo check},
and a single inference node. PACE is the exception in model scale. The
frontier systems that RRSI, HarnessBank and SAGE evolve may well show less
heterogeneous noise across configurations, and we do not claim that their
published improvements are false. We claim that, in a noisy regime, a threshold
calibrated once is a bet on the noise of configurations not yet measured, and
we quantify what that bet costs.
